\section{System and Threat Model}
\label{sec:threat}

\subsection{Parties and shared state}

Sender and receiver share an indexed image corpus $\mathcal{I}=\{x_i\}_{i=1}^{N}$, the trained descriptor and frozen codebook, and independent cryptographic subkeys for payload encryption, cluster mapping, and cover selection. The sender transmits existing indexed images unchanged. Session metadata records the authenticated sequence number, codebook size, padding information, and cover count. The implementation seals metadata with AES-GCM and binds it to sender/receiver identities.

The benchmark uses deterministic public key/nonce derivation to make every run reproducible. A production deployment can replace this benchmark root with an operational key-management service while preserving the protocol interfaces. For a partition key $K_{\mathrm{master}}$, HMAC-SHA256 derives subkeys for encryption, mapping, and selection. The benchmark nonce combines a four-byte HMAC-derived prefix with the unique 64-bit session sequence, giving explicit nonce uniqueness for sequences $0,\ldots,29$ under a fixed benchmark key.

\subsection{Channel and warden}

The image channel applies lossy compression, additive Gaussian noise, blur, resizing, central cropping, and small rotations. Calibration and holdout transformation values are disjoint. The receiver decodes directly from the received image stream under the declared channel family.

For detectability, we use a strong matched-index warden: for each authenticated session, images selected by that session form the positive class and the remaining images in the same frozen index form the negative class. Training and test image identities are disjoint. This measures selection leakage relative to the shared index distribution. The present scope focuses on image-content discrimination; sequence-level traffic features, temporal metadata, semantic profiling, and cross-platform behavioral traces define complementary future threat models.

\subsection{Security goals and claim boundary}

NARCIS separates cryptographic protection from empirical carrier-selection security. AES-GCM provides payload confidentiality/integrity and metadata authentication under its standard AEAD assumptions and nonce discipline \cite{nistgcm}; the replay guard rejects sequence numbers at or below the highest accepted value for a sender. Reed--Solomon coding \cite{reed1960polynomial} corrects residual byte errors within its decoding budget. Selection leakage is measured independently through the detector battery, while traffic-analysis resistance and broader steganographic secrecy belong to a larger deployment threat model.

\section{NARCIS Protocol}
\label{sec:method}

\begin{figure}[t]
\centering
\begin{tikzpicture}[
  node distance=5mm and 7mm,
  box/.style={draw,rounded corners,align=center,minimum height=8mm,text width=28mm,font=\small},
  arrow/.style={-{Latex[length=2mm]},thick}
]
\node[box] (train) {Learn normalized\\image descriptor};
\node[box,right=of train] (project) {Calibration-locked\\projection + $K=8$ quantiles};
\node[box,right=of project] (bank) {Five-cover group bank\\+ failure signatures};
\node[box,below=of project] (crypto) {AES-GCM + framing\\+ RS128};
\node[box,right=of crypto] (map) {Session-balanced\\keyed Gray map};
\node[box,right=of map] (select) {Signature-balanced\\keyed group selection};
\node[box,below=of map] (channel) {Image channel};
\node[box,left=of channel] (decode) {Majority labels\\+ RS decode};
\node[box,left=of decode] (auth) {Authenticated\\plaintext decision};
\draw[arrow] (train)--(project);
\draw[arrow] (project)--(bank);
\draw[arrow] (crypto)--(map);
\draw[arrow] (map)--(select);
\draw[arrow] (bank)--(select);
\draw[arrow] (select)--(channel);
\draw[arrow] (channel)--(decode);
\draw[arrow] (decode)--(auth);
\end{tikzpicture}
\caption{NARCIS separates offline channel-aware index construction (top) from online protected transmission and recovery (bottom). The external-validation operating point is frozen before Caltech-101 outcomes are inspected.}
\Description{A block diagram showing descriptor learning, projection and quantile codebook construction, five-cover group banks, AES-GCM and Reed-Solomon framing, keyed Gray mapping, keyed group selection, the image channel, majority decoding, and final authentication.}
\label{fig:pipeline}
\end{figure}

\subsection{Invariant descriptor}

A convolutional encoder $f_\theta$ maps an image to an $L_2$-normalized descriptor $z=f_\theta(x)\in\mathbb{R}^{64}$. The frozen encoder contains four convolutional blocks with batch normalization and SiLU activations, adaptive average pooling, and a two-layer projector. Training uses paired channel augmentations and variance--invariance--covariance regularization in the spirit of VICReg \cite{bardes2022vicreg}. In the external campaign, images are fit to 256$\times$256 RGB at the channel boundary and 128$\times$128 at the encoder boundary; base width is 16, batch size 16, and training lasts five epochs with AdamW ($10^{-3}$ learning rate, $10^{-5}$ weight decay).

\subsection{Calibration-locked codebook and qualification}

For a candidate unit direction $u$, each indexed image receives scalar score
\begin{equation}
  s_i = u^\top f_\theta(x_i).
\end{equation}
A balanced one-dimensional quantile codebook partitions the clean scores into $K=8$ labels. Candidate projections comprise the first 16 clean-descriptor principal directions and 32 deterministic random directions. Projection selection uses the twelve calibration transformations: JPEG qualities 80 and 50; Gaussian noise standard deviations 5 and 12; Gaussian blur radii 0.8 and 1.5; resize scales 0.75 and 0.50; central crops 5\% and 10\%; and rotations 3 and 7 degrees.

Let $y_i$ be the clean label and $y_{i,a}$ the label under calibration attack $a$. Define correctness
\begin{equation}
 c_{i,a}=\mathbb{1}[y_{i,a}=y_i].
\end{equation}
The final protocol exploits complementary error patterns across covers. Projection candidates are ranked by analytical lower bounds on unavoidable majority-failing groups, followed by stable-cover and distribution-balance diagnostics. Holdout transformations, semantic category labels, detector outcomes, and final plaintext-recovery outcomes remain reserved for later validation stages.

\subsection{Complementary group bank}

Within each codebook label, all covers of the 7,000-image index carrying that label are partitioned into disjoint groups of five. A group $G$ fails calibration attack $a$ when fewer than three members retain their clean label:
\begin{equation}
 F_a(G)=\mathbb{1}\!\left[\sum_{i\in G} c_{i,a}<3\right].
\end{equation}
The vector $(F_a(G))_a$ is stored as a calibration majority-failure signature. Group-bank construction minimizes majority-failing group/attack cells using deterministic restarts and swap proposals, and compares the achieved count against an analytical lower bound. Covers with complementary individual failure modes can therefore form a majority-correct communication unit.

During emission, groups of a requested label are partitioned by failure signature. If $p_s$ is the natural proportion of signature $s$ in the complete label-specific bank, the scheduler chooses each next available signature by its current proportional deficit. Secret-dependent HMAC ranking supplies deterministic ordering and tie breaking inside that rule. The emitted prefix consequently tracks the complete bank's signature mixture while preserving keyed group choice.

\subsection{Authenticated coding and exact session balancing}

The plaintext is encrypted with AES-GCM using associated data that include the authenticated sequence number. The protected envelope is framed and encoded with 128 Reed--Solomon parity bytes, then divided into $b=\log_2K=3$-bit symbols.

NARCIS uses an exact cyclic keyed Gray mapping to distribute each coded symbol across visual clusters over sessions. Let $g(p)=p\oplus(p\!\gg\!1)$ be binary-reflected Gray order. A mapping subkey fixes a secret base $b_0\in\{0,\ldots,K-1\}$ and one orientation $r\in\{-1,+1\}$. For authenticated session sequence $q$,
\begin{equation}
 \mathrm{shift}_q=(b_0+q)\bmod K,
\end{equation}
\begin{equation}
 \pi_q(g(p))=(\mathrm{shift}_q+r p)\bmod K.
 \label{eq:mapping}
\end{equation}
For any fixed coded symbol, Eq.~\eqref{eq:mapping} visits every cluster exactly once over any $K$ consecutive sequence values. This algebraic property supplies exact long-run balancing independently of detector outcomes.

\begin{figure}[t]
\centering
\begin{tikzpicture}[font=\scriptsize]
\matrix (m) [matrix of nodes,
  nodes={draw,minimum width=7mm,minimum height=5.5mm,anchor=center},
  row sep=-\pgflinewidth,column sep=-\pgflinewidth] {
  $q$ & $g_0$ & $g_1$ & $g_2$ & $g_3$ & $g_4$ & $g_5$ & $g_6$ & $g_7$ \\
  0 & 0 & 1 & 2 & 3 & 4 & 5 & 6 & 7 \\
  1 & 1 & 2 & 3 & 4 & 5 & 6 & 7 & 0 \\
  2 & 2 & 3 & 4 & 5 & 6 & 7 & 0 & 1 \\
  3 & 3 & 4 & 5 & 6 & 7 & 0 & 1 & 2 \\
  4 & 4 & 5 & 6 & 7 & 0 & 1 & 2 & 3 \\
  5 & 5 & 6 & 7 & 0 & 1 & 2 & 3 & 4 \\
  6 & 6 & 7 & 0 & 1 & 2 & 3 & 4 & 5 \\
  7 & 7 & 0 & 1 & 2 & 3 & 4 & 5 & 6 \\
};
\end{tikzpicture}
\caption{Illustrative normalized $K=8$ cyclic schedule for $b_0=0$ and positive orientation. Secret base and orientation permute the displayed labels while preserving the property that each coded symbol visits all eight clusters once per complete cycle.}
\Description{A nine-column matrix showing sessions zero through seven and a cyclic rotation of cluster labels zero through seven for eight Gray-ordered coded symbols.}
\label{fig:cyclic-map}
\end{figure}

For each required mapped label, the sender chooses one five-cover group using the signature-balanced keyed scheduler and transmits the five unchanged images. Each cover appears at most once inside one transmission. The receiver re-embeds each image, recovers five labels, takes their strict majority, inverts $\pi_q$, reconstructs the coded bitstream, applies Reed--Solomon decoding, and accepts plaintext after framing and AES-GCM verification.

\subsection{Finite-index feasibility}

For protected envelope $c$, let $d_k(c)$ denote the number of group occurrences required for label $k$ after coding and session mapping, and $a_k$ the number of available groups in label bank $k$. The session is feasible when
\begin{equation}
 d_k(c)\le a_k \quad \text{for every }k.
 \label{eq:feasible}
\end{equation}
Nominal alphabet capacity $\log_2 K$ gives coded bits per label, while Eq.~\eqref{eq:feasible} tests complete-message realizability under finite group availability and within-session unique-cover usage.

\subsection{Protocol properties and component roles}

Two exact properties hold before experimental evaluation. First, for any fixed coded symbol $v$ and any starting sequence $q_0$, the multiset $\{\pi_{q_0+t}(v):0\le t<K\}$ contains every cluster label exactly once because the additive shift traverses all residues modulo $K$ while Gray ordering and orientation permute fixed symbol positions. Second, under unique-cover usage and a preconstructed group bank, Eq.~\eqref{eq:feasible} is necessary and sufficient for emitting the complete coded session: each occurrence consumes one distinct group of its mapped label, and any feasible demand vector can be assigned distinct available groups label by label.

\begin{table}[t]
\caption{Role of the principal components and the evidence that supports each role.}
\label{tab:roles}
\small
\begin{tabularx}{\textwidth}{@{}lXX@{}}
\toprule
Component & Protocol role & Evidence/guarantee \\
\midrule
Calibration-locked projection & Define an externally fixed $K=8$ label space & Selection uses calibration attacks; holdouts remain reserved for validation \\
Complementary five-cover bank & Convert heterogeneous per-cover failures into majority decisions & Group failure signatures and 1,800/1,800 external calibration recoveries \\
RS128 & Correct residual coded-byte errors & Up to 37 corrections in calibration and 64 on holdout crop \\
Cyclic keyed Gray map & Balance coded-symbol-to-cluster assignment & Exact $K$-session algebraic balancing property \\
Signature-balanced scheduler & Preserve calibration-failure-signature mixture during keyed selection & Proportional-deficit scheduling followed by independent detector audit \\
AES-GCM + sequence/replay state & Authenticate accepted plaintext and session metadata & Standard AEAD semantics, explicit nonce discipline, stateful replay rejection \\
\bottomrule
\end{tabularx}
\end{table}

Table~\ref{tab:roles} isolates component roles and evidence types. The manuscript evaluates the integrated frozen protocol, while component-removal causal studies remain a natural extension for future work.
